\begin{filecontents*}[overwrite]{verify_dfinite_curve.py}
"""Exact tests of the differential matrix order-degree curve.

Requires SymPy. The proof is in the accompanying manuscript; these finite
computations check the implementation, polynomial identities, and kernel counts.
"""
from pathlib import Path
import json
import time
import sympy as S
from sympy.polys.matrices import DomainMatrix

x, y = S.symbols('x y')

def poly(v):
    return S.Poly(v, x, y, domain=S.QQ)

def verify(name, A, B, D, r=None, d=None):
    start = time.monotonic()
    m = A.rows
    D = S.sympify(D)
    Cx, Cy = A/D, B/D
    curvature = Cx.diff(y)-Cy.diff(x)+Cx*Cy-Cy*Cx
    assert all(S.cancel(v) == 0 for v in curvature)
    entries = [D]+list(A)+list(B)
    t = max([0]+[int(poly(v).degree(x)) for v in entries if v != 0])
    s = max([0]+[int(poly(v).degree(y)) for v in entries if v != 0])
    if r is None:
        r = m*(s+1)
    if d is None:
        d = m*t*((2*r-1)*s+r)//(r+1-m*(s+1))
    surplus = (r+1-m*(s+1))*(d+1)-m*t*((2*r-1)*s+r)
    assert surplus > 0
    a, b = d+(r-1)*t, (r-1)*s
    np = (r+1)*(d+1)
    nu = m*(a+1)*(b+1)
    E = m*(d+r*t+1)*(r*s+1)
    N = [S.Matrix([[1]+[0]*(m-1)])]
    for i in range(r):
        V = D*N[-1].diff(x)+N[-1]*A-i*S.diff(D,x)*N[-1]
        N.append(V.applyfunc(S.expand))
    T = [(N[i]*D**(r-i)).applyfunc(S.expand) for i in range(r+1)]
    columns = []
    for i in range(r+1):
        for j in range(d+1):
            columns.append([poly(v*x**j) for v in T[i]])
    for ell in range(m):
        for j in range(a+1):
            for k in range(b+1):
                u = S.Matrix([[0]*m])
                u[ell] = x**j*y**k
                V = -D*u.diff(y)-u*B+(r-1)*S.diff(D,y)*u
                columns.append([poly(v) for v in V])
    assert len(columns) == np+nu
    rows = {}
    for col, components in enumerate(columns):
        for ell, p in enumerate(components):
            for (ix,iy), val in p.terms():
                assert ix <= d+r*t and iy <= r*s
                if val:
                    row = ell*(d+r*t+1)*(r*s+1)+ix*(r*s+1)+iy
                    rows.setdefault(row,{})[col] = S.QQ.convert(val)
    dm = DomainMatrix(rows,(E,np+nu),S.QQ)
    ns = dm.nullspace().to_Matrix()
    good = [list(ns.row(i)) for i in range(ns.rows)
            if any(ns[i,j] for j in range(np))]
    assert good
    v = min(good,key=lambda w:sum(bool(q) for q in w))
    den = S.ilcm(*[S.denom(z) for z in v])
    v = [int(den*z) for z in v]
    content = S.igcd(*v)
    v = [int(z//content) for z in v]
    ps = [sum(v[i*(d+1)+j]*x**j for j in range(d+1)) for i in range(r+1)]
    U = S.Matrix([[sum(v[np+(ell*(a+1)+j)*(b+1)+k]*x**j*y**k
                       for j in range(a+1) for k in range(b+1))
                    for ell in range(m)]])
    residual = sum((ps[i]*T[i] for i in range(r+1)),S.zeros(1,m))
    residual -= D*U.diff(y)+U*B-(r-1)*S.diff(D,y)*U
    assert all(S.expand(v) == 0 for v in residual)
    # Independently verify the un-cleared rational connection identity.
    R = U/D**(r-1)
    raw = sum((ps[i]*N[i]/D**i for i in range(r+1)),S.zeros(1,m))
    raw -= R.diff(y)+R*Cy
    assert all(S.cancel(v) == 0 for v in raw)
    ur = {i:{j-np:c for j,c in row.items() if j >= np} for i,row in rows.items()}
    ur = {i:row for i,row in ur.items() if row}
    udm = DomainMatrix(ur,(E,nu),S.QQ)
    pure_kernel = nu-udm.rank()
    assert pure_kernel <= m*(a+1)
    norm = lambda p: sum(abs(c) for c in poly(p).coeffs())
    alpha = max(sum(norm(A[i,j]) for j in range(m)) for i in range(m))
    beta = max(sum(norm(B[i,j]) for j in range(m)) for i in range(m))
    delta = norm(D)
    HP = max(delta**(r-i)*S.prod(alpha+2*j*t*delta for j in range(i))
             for i in range(r+1))
    HU = max(1,beta+2*(r-1)*s*delta)
    assert all(sum(sum(abs(v) for v in p.coeffs()) for p in col) <= HP
               for col in columns[:np])
    assert all(sum(sum(abs(v) for v in p.coeffs()) for p in col) <= HU
               for col in columns[np:])
    K = min(max(HP,HU)**E, HP**np*HU**nu)
    assert max(abs(z) for z in v) <= K
    result = dict(name=name,m=m,t=t,s=s,r=r,d=d,rows=E,columns=np+nu,
                  positive_surplus=surplus,kernel_dimension=ns.rows,
                  pure_certificate_kernel=pure_kernel,
                  pure_kernel_bound=m*(a+1),
                  actual_order=max(i for i,p in enumerate(ps) if p != 0),
                  actual_degree=max(int(S.degree(p,x)) for p in ps if p != 0),
                  P=[str(S.factor(p)) for p in ps],
                  certificate_row=[str(S.factor(v)) for v in R],
                  height_bound_bits=int(K).bit_length(),
                  seconds=round(time.monotonic()-start,3),verified=True)
    print(json.dumps(result),flush=True)
    return result

if __name__ == '__main__':
    Z = S.zeros(1)
    M = S.Matrix([[0,1,0],[0,0,1],[1,1,0]])
    cases = [
        ('constant',Z,Z,1),
        ('sine translate',S.Matrix([[0,1],[-1,0]]),S.Matrix([[0,1],[-1,0]]),1),
        ('rank-three polynomial exponential translate',x*M,M,1),
        ('Airy translate',S.Matrix([[0,1],[x+y,0]]),S.Matrix([[0,1],[x+y,0]]),1),
        ('exponential times Airy',y*S.eye(2),S.Matrix([[x,1],[y,x]]),1),
        ('Bessel translate',S.Matrix([[0,x+y],[-(x+y),-1]]),S.Matrix([[0,x+y],[-(x+y),-1]]),x+y),
        ('rank-two nilpotent flat module',S.zeros(2),S.Matrix([[0,1],[0,0]]),1),
    ]
    result = [verify(*case) for case in cases]
    Path(__file__).with_name('dfinite_curve_checks.json').write_text(json.dumps(result,indent=2)+'\n')
\end{filecontents*}
\begin{filecontents*}[overwrite]{verify_special_function_certificates.py}
"""Check exact rank-two special-function certificates in their ODE modules."""
import sympy as S
from pathlib import Path
import json
x,y = S.symbols('x y')

def check(name,Cx,Cy,coefficients,certificate):
    curvature = Cx.diff(y)-Cy.diff(x)+Cx*Cy-Cy*Cx
    assert all(S.cancel(v)==0 for v in curvature)
    v = S.Matrix([[1,0]])
    lhs = S.zeros(1,2)
    for p in coefficients:
        lhs += p*v
        v = (v.diff(x)+v*Cx).applyfunc(S.cancel)
    residual = lhs-certificate.diff(y)-certificate*Cy
    assert all(S.cancel(v)==0 for v in residual)
    return dict(name=name,coefficients=list(map(str,coefficients)),
                certificate=list(map(str,certificate)),verified=True)

if __name__=='__main__':
    z=x+y*y
    airy=S.Matrix([[0,1],[z,0]])
    bessel=S.Matrix([[0,1],[-1,-1/z]])
    data=[
        check('Airy quadratic pullback',airy,2*y*airy,[-1,-2*x,0,2],S.Matrix([[y,0]])),
        check('Bessel quadratic pullback',bessel,2*y*bessel,[1,2*x,3,2*x],S.Matrix([[0,y/z]])),
        check('Airy product pullback',S.Matrix([[0,y],[x*y*y,0]]),
              S.Matrix([[0,x],[x*x*y,0]]),[1,x],S.Matrix([[y,0]]))
    ]
    print(json.dumps(data,indent=2))
    Path(__file__).with_name('special_function_checks.json').write_text(json.dumps(data,indent=2)+'\n')
\end{filecontents*}
\documentclass[11pt,a4paper]{article}
\usepackage[margin=25mm]{geometry}
\usepackage[T1]{fontenc}
\usepackage{lmodern,amsmath,amssymb,amsthm,mathtools,booktabs,array}
\usepackage{xcolor,microtype,fancyhdr}
\usepackage[colorlinks=true,linkcolor=blue!45!black,urlcolor=blue!45!black,citecolor=blue!45!black]{hyperref}
\hypersetup{pdftitle={Explicit Order-Degree Curves for Differential D-Finite Telescoping},pdfauthor={Keio}}
\newtheorem{theorem}{Theorem}[section]
\newtheorem{lemma}[theorem]{Lemma}
\newtheorem{proposition}[theorem]{Proposition}
\newtheorem{corollary}[theorem]{Corollary}
\theoremstyle{remark}\newtheorem{remark}[theorem]{Remark}
\setlength{\emergencystretch}{3em}
\setlength{\headheight}{14pt}
\pagestyle{fancy}\fancyhf{}
\fancyhead[L]{\small Keio}
\fancyhead[R]{\small Order--degree curves for differential telescoping}
\fancyfoot[C]{\thepage}
\title{Explicit Order--Degree Curves for\\Differential D-Finite Telescoping\\[5pt]
\large A connection-matrix refinement addressing Chen--Kauers Problem 3}
\author{Keio\thanks{Research development, computational checks, and manuscript preparation with ChatGPT.}}
\date{2 November 2026}
\begin{document}
\maketitle
\begin{abstract}
For a bivariate differential D-finite input presented by rank-$m$ rational
connection matrices $A/D$ and $B/D$, we give an explicit sufficient
order--degree curve using only the maximum bidegrees $(t,s)$ of $A,B,D$.
For every integer $r\ge m(s+1)$, a telescoper of exact order $r$ exists with
degree at most
\[
 \left\lfloor\frac{mt[(2r-1)s+r]}{r+1-m(s+1)}\right\rfloor.
\]
A polynomial certificate support bound and an integer coefficient bound
are obtained simultaneously. The argument counts polynomial unknowns
in the generalized Apagodu--Zeilberger ansatz and controls the
certificate-only kernel by evaluation at an ordinary point. It imposes
no Fuchsian or regular-singularity assumption. The result addresses
the differential setting of Chen--Kauers Problem~3; it does not establish
a first resolution, a new foundational principle, or an extension to
arbitrary discrete and mixed Ore systems. Seven exact ansatz computations
and three special-function certificate identities accompany the proof.
Their verification scripts are embedded in the single TeX source.
\end{abstract}

\section{Question and contribution}
Chen and Kauers' original arXiv Problem~3 asks for an order--degree curve
for general D-finite functions~\cite{CK2016}. Their definition permits
a sufficient bound above an order threshold; it does not demand the
true minimum order or a sharp boundary. This note supplies a concrete
formula in the full bivariate \emph{differential} setting. The problem's
surrounding discussion also includes summation, which is not covered here.

The finite input used in the theorem consists of derivative multiplication
matrices. It is available from any finite-rank annihilating ideal by
standard normal forms. No telescoper or reduced integration quotient
must be computed to evaluate the bound. This separates the result from
a bound that uses the output operator itself as an input parameter.

\paragraph{Attribution and status.}
The vector ansatz and certificate-kernel subtraction already occur in
Chen--Kauers--Koutschan~\cite[Theorem~6]{CKK2014}; the present calculation
tracks both polynomial degrees explicitly. Kauers' book discusses
order--degree bounds in its broader treatment of creative
telescoping~\cite[Section~5.5]{Kauers2023}. Van der Hoeven's reduction
framework gives polynomial complexity and degree control for nonminimal
differential telescopers in its representation~\cite[Section~9]{vdH2018}.
Consequently, general degree control is not itself a new discovery here.
The particular direct formula below was not identified in the sources
examined, but that observation does not establish historical priority.
This manuscript presents a proved quantitative refinement, not a
certified first solution of a currently open problem.

\paragraph{What is proved.}
The proof covers every finite-rank differential presentation over a
characteristic-zero field, including irregular singularities. It also
provides explicit support bounds for a certificate and, over the
rationals, an arithmetic extension of the rank-one construction.
The proof and finite checks have been reviewed during preparation;
they have not undergone external peer review or proof-assistant
formalization.
\section{An input-level order--degree curve for differential D-finite functions}
\label{gd:section}

Chen and Kauers' original Problem~3 asks for an order--degree curve for
general D-finite functions. We establish such a curve for every bivariate
\emph{differential} D-finite input, expressed in its rational connection
matrices. These matrices are computed from the input annihilating ideal;
no telescoper or reduced telescoping quotient is assumed available.
The construction refines the generalized Apagodu--Zeilberger polynomial
ansatz of Chen, Kauers, and Koutschan~\cite{CKK2014} by tracking both polynomial
degrees. It is not asserted to be a historically new methodological
principle. Discrete and mixed Ore variants are outside the theorem.

\subsection{Input representation}
\label{gd:representation}

Let $\mathbb K$ be a characteristic-zero field, and put
$\mathcal R=\mathbb K(x,y)\langle\partial_x,\partial_y\rangle$.
Suppose first that a nonzero D-finite function $f$ is given together with
a basis $F=(F_1,\ldots,F_m)^T$ of its derivative module over
$\mathbb K(x,y)$, with $F_1=f$. Such a basis exists: the derivative
module is finite-dimensional and contains the nonzero element $f$.
Every $F_j$ is an operator in $\mathcal R$ applied to $f$, so rational
row combinations of $F$ are permitted certificate functions. Clearing
denominators gives
\begin{equation}\label{gd:connection}
 \partial_xF=\frac AD F,\qquad
 \partial_yF=\frac BD F,
 \qquad A,B\in\mathbb K[x,y]^{m\times m},\quad
 0\ne D\in\mathbb K[x,y].
\end{equation}

The same representation is effective from a proper D-finite left ideal
$I\subseteq\mathcal R$ of finite rank $m$. Compute a left Gr\"obner
basis and enumerate its standard differential monomials
$b_1=1,b_2,\ldots,b_m$. Normal forms of
$\partial_x b_j$ and $\partial_y b_j$ give the two multiplication
matrices. A common denominator supplies~\eqref{gd:connection} in the
formal quotient module $\mathcal R/I$, with $F_j=[b_j]$ and
$F_1=[1]$. If $I$ is merely contained in the annihilator of a particular
function, the evaluated $b_j(f)$ may be dependent; this causes no problem,
because the formal identity modulo $I$ implies the identity on $f$.
If $I=\mathcal R$, the represented function is zero and $P=1$ is trivial.

The parameters below are explicit in this computed connection
presentation. We do not claim a conversion bound solely in terms of the
degrees of an arbitrary unprocessed PDE generating set. All presentation
steps precede the telescoping calculation and use only the supplied
D-finite ideal. Over an effectively computable coefficient field, they
are ordinary finite normal-form and rational-function computations.

\subsection{The order--degree theorem}

\begin{theorem}\label{gd:curve}
Assume~\eqref{gd:connection}, with $m\ge1$, and let $t,s\ge0$ be integer
upper bounds on the $x$- and $y$-degrees of $D$ and every entry of $A,B$.
Let $r,d$ be nonnegative integers such that
\begin{equation}\label{gd:condition}
 r+1>m(s+1),\qquad
 \bigl[r+1-m(s+1)\bigr](d+1)
   >mt\bigl[(2r-1)s+r\bigr].
\end{equation}
Set
\begin{equation}\label{gd:boxes}
 a=d+(r-1)t,\qquad b=(r-1)s.
\end{equation}
There exist polynomials $p_0,\ldots,p_r\in\mathbb K[x]$, not all zero,
and a row $U\in\mathbb K[x,y]^{1\times m}$ such that
\begin{align}
 &\deg_xp_i\le d,\qquad
   \deg_xU_j\le a,\qquad \deg_yU_j\le b,\label{gd:degrees}\\
 &\sum_{i=0}^r p_i(x)\partial_x^iF_1
    =\partial_y\left(D^{1-r}UF\right).
       \label{gd:telescoping}
\end{align}
In the formal quotient presentation, this states
$P-\partial_yQ\in I$, where
\[
 P=\sum_{i=0}^rp_i(x)\partial_x^i\ne0,\qquad
 Q=D^{1-r}\sum_{j=1}^m U_jb_j\in\mathcal R.
\]
Consequently, for every integer $r\ge m(s+1)$ there is a telescoper
of \emph{exact} order $r$ and degree at most $\lfloor c(r)\rfloor$, where
\begin{equation}\label{gd:hyperbola}
 \boxed{\displaystyle
 c(r)=\frac{mt\bigl[(2r-1)s+r\bigr]}{r+1-m(s+1)}.}
\end{equation}
\end{theorem}

\begin{proof}
Let $v_0=(1,0,\ldots,0)$ and recursively define polynomial rows
\begin{equation}\label{gd:derivative}
 v_{i+1}=D\partial_xv_i+v_iA-i(\partial_xD)v_i.
\end{equation}
By induction, $\partial_x^iF_1=D^{-i}v_iF$ and every entry of $v_i$
has bidegree at most $(it,is)$. Put $w_i=v_iD^{r-i}$, whose entries
have bidegree at most $(rt,rs)$.

Make an ansatz with the degree bounds~\eqref{gd:degrees}.
The coordinate identity sufficient for~\eqref{gd:telescoping} is
\begin{equation}\label{gd:coefficient}
 \sum_{i=0}^r p_iw_i
  =D\partial_yU+UB-(r-1)(\partial_yD)U.
\end{equation}
Every coordinate has bidegree at most $(d+rt,rs)$. With
\[
 n_P=(r+1)(d+1),\qquad
 n_U=(a+1)(b+1),\qquad
 E=(d+rt+1)(rs+1),
\]
coefficient comparison yields a homogeneous linear system over
$\mathbb K$ with at most $mE$ equations and $n_P+mn_U$ unknowns.

We exclude solutions supported only on the certificate coordinates.
Such a row $U$ satisfies
\begin{equation}\label{gd:homogeneous}
 \partial_yU
   =U\frac{(r-1)(\partial_yD)I_m-B}{D}.
\end{equation}
Choose $y_0\in\{0,1,\ldots,s\}$ with $D(x,y_0)\ne0$ in
$\mathbb K[x]$; at least one exists because $\deg_yD\le s$.
Over $\mathbb K(x)$ the coefficient matrix in~\eqref{gd:homogeneous}
is regular at $y_0$. Formal power-series substitution around $y_0$
shows uniqueness for a prescribed value $U(x,y_0)$: the coefficient
of $(y-y_0)^{j+1}$ is determined by earlier coefficients after division
by the nonzero integer $j+1$. Therefore evaluation at $y_0$ injects
the bounded polynomial solution space into
$(\mathbb K[x]_{\le a})^m$. Its dimension over $\mathbb K$ is at most
$m(a+1)$.

The full solution space has dimension at least $n_P+mn_U-mE$, and
direct expansion gives
\begin{align*}
 n_P+mn_U-mE-m(a+1)
 &=\bigl[r+1-m(s+1)\bigr](d+1)\\
 &\hspace{1em}-mt\bigl[(2r-1)s+r\bigr]>0.
\end{align*}
It consequently contains a vector with some $p_i\ne0$, proving
the first assertion. Notice that only coordinate matrices enter the
argument; no analytic regularity of the solution functions is required.

For the curve, choose $d=\lfloor c(r)\rfloor$. Then
$d+1>c(r)$, so~\eqref{gd:condition} holds. The constructed operator
may have order $k<r$. In that case replace $P$ by
$\partial_x^{r-k}P$ and $Q$ by $\partial_x^{r-k}Q$.
Commutation of the two derivations preserves the certificate identity.
The new operator has order exactly $r$, and its coefficient degrees
do not increase, since left multiplication by $\partial_x$ only
differentiates existing coefficients or shifts their operator indices.
This proves~\eqref{gd:hyperbola} with the exact-order convention.
\end{proof}

\begin{proposition}\label{gd:balanced}
Two uniform admissible pairs in Theorem~\ref{gd:curve} are
\begin{align*}
 r&=m(s+1),&
 d&=mt\bigl[m(s+1)(2s+1)-s\bigr],\\
 r&=2m(s+1)-1,&
 d&=t(4ms+2m-1).
\end{align*}
The second gives order $O(m(s+1))$ and degree $O(mt(s+1))$.
The curve~\eqref{gd:hyperbola} has horizontal asymptote $mt(2s+1)$.
\end{proposition}
\begin{proof}
For the first pair the surplus in the strict inequality
of~\eqref{gd:condition} is one. For the second it is
$m(s+1+2ts)>0$. The asymptote follows by dividing the numerator and
denominator of~\eqref{gd:hyperbola} by $r$.
\end{proof}

\begin{remark}\label{gd:sharpness}
The formula gives a sufficient region and does not identify the true
minimum order, optimal degree, or a generic sharp curve. Neither
$t=0$ nor $s=0$ needs exclusion. The exact-order adjustment in the
proof preserves the degree bound, but it is not asserted to preserve
the original certificate box in~\eqref{gd:degrees}.
\end{remark}

\subsection{An arithmetic extension over the rationals}

\begin{proposition}\label{gd:height}
Suppose $A,B,D$ have integer coefficients. Define
\begin{align*}
 \alpha&=\max_i\sum_j\|A_{ij}\|_1,&
 \beta&=\max_i\sum_j\|B_{ij}\|_1,&
 \delta&=\|D\|_1,\\
 H_P&=\max_{0\le i\le r}
  \left\{\delta^{r-i}\prod_{j=0}^{i-1}(\alpha+2jt\delta)\right\},\\
 H_U&=\max\{1,\beta+2(r-1)s\delta\},&
 H&=\max\{H_P,H_U\}.
\end{align*}
Here $\|\cdot\|_1$ denotes the coefficient $\ell_1$ norm and an empty
product equals one. With $n_P,n_U,E$ as in the proof of
Theorem~\ref{gd:curve}, one may choose integer $p_i,U_j$ in
\eqref{gd:telescoping}, with some $p_i\ne0$, such that every coefficient
has absolute value at most
\begin{equation}\label{gd:heightbound}
 \min\{H^{mE},H_P^{n_P}H_U^{mn_U}\}.
\end{equation}
This assertion bounds the operator of order at most $r$ before any
exact-order adjustment. For $m,t,s\ge1$, the balanced pair in
Proposition~\ref{gd:balanced} gives
\[
 O\!\left(m^3ts^3\,[h+\log(mts+2)]\right)
\]
bits per coefficient, where $h=\log_2\max\{1,\alpha,\beta,\delta\}$.
\end{proposition}

\begin{proof}
For polynomial rows use the sum of the coefficient $\ell_1$ norms of
their entries. Then $\|vA\|_1\le\alpha\|v\|_1$, and
recurrence~\eqref{gd:derivative} gives
\[
 \|v_{i+1}\|_1\le(\alpha+2it\delta)\|v_i\|_1.
\]
Every telescoper column in the coefficient system has $\ell_1$ norm
at most $H_P$. A certificate column involves one row of $B$ and one
scalar derivative contribution; its norm is at most
$\beta+2(r-1)s\delta\le H_U$.

Choose independent columns and a nonsingular maximal-rank square
submatrix. Setting one free coordinate to its determinant and solving
for the pivot coordinates by the adjugate produces integer kernel
vectors whose coordinates are maximal-rank minors. They span the
rational kernel. Since that kernel has a nonzero telescoper projection,
one such vector has a nonzero telescoper projection. Hadamard's
inequality bounds a minor by $H^{mE}$ and, by counting the two types
of columns, by $H_P^{n_P}H_U^{mn_U}$. If the matrix has rank zero,
use a unit vector in a telescoper coordinate instead.

For the balanced choice,
$n_P=O(m^2ts^2)$ and $mn_U=O(m^3ts^3)$, while
$\log H_P=O(ms[h+\log(mts+2)])$ and
$\log H_U=O(h+\log(ms+2))$.
Taking logarithms of the second bound in~\eqref{gd:heightbound}
proves the assertion.
\end{proof}

\subsection{Relation to the original problem and prior work}

Theorem~\ref{gd:curve} covers the full class of finite-rank
\emph{differential} D-finite presentations, rather than a selected
function family. Its constants can be calculated directly from
derivative multiplication matrices before telescoping. This supplies
an explicit sufficient order--degree curve for the integration setting
of original Problem~3.

The generalized Apagodu--Zeilberger theorem of Chen, Kauers, and
Koutschan~\cite{CKK2014} already uses these multiplication matrices and bounds the
dimension of the certificate-only kernel. Its ansatz has coefficients
in $\mathbb K(x)$ and gives an order bound. The calculation here bounds
the unknown coefficients polynomially in $x$, counts both degrees,
and yields~\eqref{gd:hyperbola}. Thus the formula is an explicit
bidegree refinement of an established method. Independent verification
of this derivation does not establish absence of an identical prior
bound or a first-resolution claim. Since the original problem's
surrounding discussion also includes summation, the result should not
be described as resolving all shift, $q$-shift, or mixed variants.

\section{Construction and worked identities}
\subsection{A finite algorithm}
An effective implementation follows the proof directly.
\begin{enumerate}
\item Compute $A,B,D$ from a derivative basis or an annihilating ideal,
and read off $m,t,s$.
\item Choose $r\ge m(s+1)$ and $d=\lfloor c(r)\rfloor$, or choose the
balanced pair of Proposition~\ref{gd:balanced}.
\item Generate $v_0,\ldots,v_r$ by~\eqref{gd:derivative} and expand the
coordinate equations~\eqref{gd:coefficient} in $x$ and $y$.
\item Compute a basis of the homogeneous coefficient kernel. Select a
vector with a nonzero telescoper projection, which the dimension
argument guarantees exists.
\item Recover $P,Q$ and verify $P-\partial_yQ$ by normal form modulo
the original ideal. Raise the order by left multiplication only if
an exact requested order is needed.
\end{enumerate}
This gives an algorithm over an effectively computable field.
The curve is measured after construction of the connection presentation;
no bound on the cost of converting arbitrary input PDE generators is
included. The proof also works with an annihilating ideal strictly
smaller than the full annihilator, though its rank and degree bounds
may then be less economical.

\subsection{Rank two: an exponential times an Airy solution}
Let $g''(y)=yg(y)$ and set
\[
 f=e^{xy}g(y),\qquad
 F=\begin{pmatrix}e^{xy}g(y)\\e^{xy}g'(y)\end{pmatrix}.
\]
The formal differential module has the compatible presentation
\[
 D=1,\qquad A=yI_2,\qquad
 B=\begin{pmatrix}x&1\\y&x\end{pmatrix},
 \qquad m=2,\ t=s=1.
\]
Since $F_2=(\partial_y-x)f$, the certificate belongs to the required
derivative module. Directly,
\begin{equation}\label{ex:airy-exponential}
 (\partial_x-x^2)f
   =\partial_y(F_2-xf)
   =\partial_y\bigl((\partial_y-2x)f\bigr).
\end{equation}
For this presentation the uniform curve is
$c(r)=(6r-2)/(r-3)=6+16/(r-3)$, for $r\ge4$.
Some guaranteed degree bounds are
\[
\begin{array}{c|rrrrrr}
 r&4&5&6&7&8&20\\\hline
 \lfloor c(r)\rfloor&22&14&11&10&9&6
\end{array}
\]
Identity~\eqref{ex:airy-exponential} is much smaller than these sufficient
bounds. This illustrates why a valid general curve should not be
confused with the optimal curve for a particular function.

\subsection{A quadratic Airy pullback}
Put $z=x+y^2$ and let $g''(z)=zg(z)$. For
$F=(g(z),g'(z))^T$, one has
\[
 A=\begin{pmatrix}0&1\\z&0\end{pmatrix},\qquad
 B=2yA,\qquad D=1.
\]
The parameters $m=2,t=1,s=3$ are valid. Writing $f=g(z)$,
the Airy equation gives $f_{xxx}=f+zf_x$, and therefore
\begin{equation}\label{ex:airy-quadratic}
 (2\partial_x^3-2x\partial_x-1)f
 =f+2y^2f_x=\partial_y(yf).
\end{equation}
This identity has order three and degree one. The general theorem, which
does not inspect the Airy equation's special cancellations, starts at
$r=8$ and gives $d=106$ there. Equation~\eqref{ex:airy-quadratic}
fits the theorem's ansatz at that order after zero padding.

\subsection{A moving denominator in a Bessel pullback}
Let $g$ solve the Bessel equation $zg''+g'+zg=0$ and retain $z=x+y^2$.
For $F=(g(z),g'(z))^T$ a polynomial common-denominator presentation is
\[
 D=z,\qquad
 A=\begin{pmatrix}0&z\\-z&-1\end{pmatrix},\qquad B=2yA.
\]
Again $m=2,t=1,s=3$. The certificate identity is
\begin{equation}\label{ex:bessel}
 (2x\partial_x^3+3\partial_x^2+2x\partial_x+1)f
   =\partial_y\left(\frac{y}{z}\partial_x f\right).
\end{equation}
For verification, use $f_{xx}=-f-f_x/z$ and
$f_{xxx}=-f_x+f/z+2f_x/z^2$. Both sides then equal
\[
 -\frac{2y^2}{z}f+
 \left(\frac1z-\frac{4y^2}{z^2}\right)f_x.
\]
Thus the rational certificate is valid as a differential-module identity,
including its moving algebraic denominator. No assertion about evaluating
an improper integral across a certificate pole follows automatically.

\subsection{From identities to definite integrals}
For fixed endpoints $a,b$, whenever differentiation under the integral
and the endpoint evaluations are justified, $Pf=\partial_y(Qf)$ yields
\[
 P\int_a^b f(x,y)\,dy=(Qf)(x,b)-(Qf)(x,a).
\]
A homogeneous equation for the integral requires those boundary terms
to vanish or to be handled separately. The theorem is an algebraic
creative-telescoping statement; it does not impose convergence or
boundary conditions on arbitrary branches of its input functions.

\section{Exact checks and reproducibility}
The first embedded script constructs the entire coefficient system over
$\mathbb Q$, computes its nullspace, selects a nonzero telescoper
projection, and checks both the cleared polynomial identity and the
uncleared rational connection identity. It separately computes the
certificate-only kernel and checks the column norm bounds used in
Proposition~\ref{gd:height}. All seven cases below passed.
The actual kernel dimensions are computed over $\mathbb Q$; zero
equation rows are counted in the stated matrix sizes.

\begin{center}
\small
\begin{tabular}{@{}lrrrr@{}}
\toprule
Presentation & $(m,t,s)$ & $(r,d)$ & Matrix size & Pure kernel\\
\midrule
Constant & $(1,0,0)$ & $(1,0)$ & $1\times3$ & $1$\\
Sine translate & $(2,0,0)$ & $(2,0)$ & $2\times5$ & $0$\\
Rank-three translate & $(3,1,0)$ & $(3,9)$ & $39\times76$ & $0$\\
Airy translate & $(2,1,1)$ & $(4,22)$ & $270\times323$ & $0$\\
Exponential times Airy & $(2,1,1)$ & $(4,22)$ & $270\times323$ & $0$\\
Bessel translate & $(2,1,1)$ & $(4,22)$ & $270\times323$ & $0$\\
Nilpotent module & $(2,0,0)$ & $(2,0)$ & $2\times5$ & $1$\\
\bottomrule
\end{tabular}
\end{center}

The translate examples use $g(x+y)$ with the defining sine, Airy, or
order-zero Bessel equation. The rank-three example uses
$g(y+x^2/2)$ with $g^{(3)}=g'+g$. The nilpotent example is the formal cyclic
module with $\partial_xF=0$ and
$\partial_yF=\left(\begin{smallmatrix}0&1\\0&0\end{smallmatrix}\right)F$;
its rank is that of the supplied presentation, not necessarily of a
particular scalar solution.

The second script verifies~\eqref{ex:airy-quadratic},
\eqref{ex:bessel}, and the product-pullback identity
$(x\partial_x+1)g(xy)=\partial_y(yg(xy))$ in their rank-two differential
modules. All three residual vectors are exactly zero. These calculations
use symbolic rational arithmetic, not numerical values of Airy or Bessel
functions. Their role is to check examples and the implemented construction;
the general theorem is proved by the dimension argument, not by these tests.

\paragraph{Single-source use.}
The TeX source contains its bibliography and two Python programs in
\texttt{filecontents*} blocks. A normal compilation writes the programs
to the compilation directory. With Python and SymPy installed, run
\begin{verbatim}
python3 verify_dfinite_curve.py
python3 verify_special_function_certificates.py
\end{verbatim}
The programs emit exact-check results and save JSON records. Compiling
the PDF does not execute Python or require shell escape. No separate
bibliography database, figure file, or source fragment is needed.

\section{Scope of the outcome}
The established statement is an explicit sufficient curve, with a full
proof, for all bivariate differential D-finite inputs in a finite-rank
connection presentation. It extends the rank-one count to arbitrary
rank and gives certificate and arithmetic bounds. It does not claim
minimum order, optimal degree, basis invariance, or a curve for arbitrary
shift and mixed systems.

The historical problem predates this manuscript, and the differential
target has not been replaced by a specially chosen function family.
However, a new underlying principle and absence of identical earlier
solutions have not been established. Existing methods already supply
the main mechanism and broader degree-control results. Thus this note
does not meet the stronger standard of a certified new closure of a
still-open named problem. It records the precise result that has been
proved, with the attribution and limitations needed to assess it.

\begin{thebibliography}{9}
\bibitem{CK2016}
S. Chen and M. Kauers,
\emph{Some open problems related to creative telescoping},
arXiv:1609.03768, 2016, especially original Problem~3.
\url{https://arxiv.org/abs/1609.03768}.

\bibitem{CKK2014}
S. Chen, M. Kauers, and C. Koutschan,
\emph{A generalized Apagodu--Zeilberger algorithm},
Proceedings of ISSAC 2014, pp.~107--114.
\url{https://doi.org/10.1145/2608628.2608641};
\url{https://arxiv.org/abs/1402.2409}.

\bibitem{Kauers2023}
M. Kauers, \emph{D-Finite Functions}, Algorithms and Computation in
Mathematics, vol.~30, Springer, 2023, Section~5.5.
\url{https://doi.org/10.1007/978-3-031-34652-1}.

\bibitem{vdH2018}
J. van der Hoeven, \emph{Creative telescoping using reductions},
draft of 8 June 2018, especially Section~9 and Theorem~9.8.
\url{https://www.texmacs.org/joris/redtel/redtel.pdf}.
\end{thebibliography}
\end{document}
